\documentclass[fontsize=20pt]{scrartcl}
\usepackage[pdf]{optional}
\usepackage[utf8]{inputenc}
\usepackage{scrpage2} % пакет для настройки колонтитулов

\usepackage[english, russian]{babel}
\usepackage{paratype}
\usepackage[left=8mm,right=8mm,top=15mm,bottom=20mm,a4paper,landscape]{geometry}
\usepackage{graphicx}
\usepackage[usenames,dvipsnames]{color}
\usepackage{tikz}
\usepackage{color}
\usetikzlibrary{calendar,fit}
\usepackage{expl3,xparse}
\usepackage{comment}
\usetikzlibrary{shapes.geometric}
\usepackage[strict]{changepage}
%\usepackage[usenames,dvipsnames]{xcolor}
%\usepackage{afterpage}


\renewcommand*{\titlepagestyle}{empty}

% Adapted from http://tex.stackexchange.com/a/10199/4771 добавляем названия 
% дней недели
\makeatletter%
\tikzoption{day headings}{\tikzstyle{day heading}=[#1]}
\tikzstyle{day heading}=[]
\tikzstyle{day letter headings}=[
    execute before day scope={ \ifdate{day of month=1}{%
      \pgfmathsetlength{\pgf@ya}{\tikz@lib@cal@yshift}%
      \pgfmathsetlength\pgf@xa{\tikz@lib@cal@xshift}%
      \pgftransformyshift{-\pgf@ya}
      \foreach \d/\l in {0/Понедельник,1/Вторник,2/Среда,3/Четверг,%
                         4/Пятница,5/Суббота,6/Воскресенье} {
        \pgf@xa=\d\pgf@xa%
        \pgftransformxshift{\pgf@xa-\cellwidth/2}%
        \pgftransformyshift{\pgf@ya}%
        \node[above=-0.5ex,day heading]{\l};%
      }
    }{}%
  }%
]
\makeatother%
%% End




%% Adapted from pgf source
\def\pgfcalendarmonthname#1{%
  \translate{\ifcase#1\or ЯНВАРЬ\or ФЕВРАЛЬ\or МАРТ\or АПРЕЛЬ\or
    МАЙ\or ИЮНЬ\or ИЮЛЬ\or АВГУСТ\or СЕНТЯБРЬ\or ОКТЯБРЬ\or
    НОЯБРЬ\or ДЕКАБРЬ\fi}%
}


% переопределяем вид заголовка (название месяца)
\makeatletter
\tikzstyle{month label above centered}=[%
  execute before day scope={%
    \ifdate{day of month=1}{%
      {
        \pgfmathsetlength{\pgf@xa}{\tikz@lib@cal@xshift}%
        \pgf@xb=\tikz@lib@cal@width\pgf@xa%
        \advance\pgf@xb by-\pgf@xa%
        \pgf@xb=.5\pgf@xb%
        \pgftransformxshift{\pgf@xb}%
        \pgftransformxshift{-\cellwidth/2}%
        \pgfmathsetlength{\pgf@y}{\tikz@lib@cal@yshift}%
        \pgftransformyshift{0.333\pgf@y}
        \tikzmonthcode%
      }
    }{}},
  every month/.append style={anchor=base}
]
\makeatother
%% End

\ExplSyntaxOn

%% Adapted from http://tex.stackexchange.com/a/56214/4771
%%
% first of all we define the user level commands
\NewDocumentCommand{\addtext}{ m }{ \bdaycal_input_add:n { #1 } }
\NewDocumentCommand{\addtextyear}{ mm }{ \bdaycal_input_add_y:nn { #1 } { #2 } }
\NewDocumentCommand{\showtext}{ }{ \bdaycal_output_direct: }

% allocate variable:
% a sequence for global storage of the inputs;
\seq_new:N \g_bdaycal_input_seq

% store globally an input in the sequence
\cs_new:Npn \bdaycal_input_add:n #1
 {
  \seq_gput_right:Nn \g_bdaycal_input_seq { #1 }
 }

\cs_new:Npn \bdaycal_input_add_y:nn #1 #2
 {
  \seq_gput_right:Nn \g_bdaycal_input_seq { #1 ~ ( \int_to_arabic:n
    { \pgfcalendarifdateyear - #2 } ) }
 }

% how to output in direct order; we simply do a mapping function calling
% \bdaycal_print:n after incrementing the counter
\cs_new_protected:Npn \bdaycal_output_direct:
 {
  \seq_map_inline:Nn \g_bdaycal_input_seq
   {
    \bdaycal_print:n { ##1 }
   }
  \seq_gclear:N \g_bdaycal_input_seq
 }

% the printing macro; change here for adapting to your wishes
\cs_new:Npn \bdaycal_print:n #1
 {
  #1 \par
 }
%% End

\int_new:N \l_advent_Y_int
\int_new:N \l_advent_xmas_julian_day_int
\int_new:N \l_advent_xmas_week_day_int
\int_new:N \l_advent_julian_day_int

\cs_new:Nn \advent_sunday:n {

  \int_set:Nn \l_advent_Y_int { #1 }

  \pgfcalendardatetojulian { \l_advent_Y_int - 12 - 25 } {
    \l_advent_xmas_julian_day_int }

  \pgfcalendarjuliantoweekday { \l_advent_xmas_julian_day_int } {
    \l_advent_xmas_week_day_int }

  \int_set:Nn \l_advent_julian_day_int {
    \l_advent_xmas_julian_day_int - \l_advent_xmas_week_day_int - 22 }

}

\pgfkeys{/pgf/calendar/Advent/.default = 0}
\pgfkeys{/pgf/calendar/Advent/.code =
  {
    \advent_sunday:n { \pgfcalendarifdateyear }
    \int_compare:nNnT { \pgfcalendarifdatejulian }
    =  {\l_advent_julian_day_int + #1}
    { \pgfcalendarmatchestrue }
  }
}

\definecolor{roed}{rgb}{0.937254901961,0.16862745098,0.176470588235}
\definecolor{navy}{rgb}{0.000000,0.000000,0.501961}
\definecolor{dimgray}{rgb}{0.411765,0.411765,0.411765}
\definecolor{lightgray}{rgb}{0.827451,0.827451,0.827451}
%%% выводим номер недели

\pgfkeys{/tikz/week~number/.code =
  {
    \week_number:nnn { 
      \pgfcalendarifdateyear } { 
      \pgfcalendarifdatemonth } { 
      \pgfcalendarifdateday }
    \addtext{ \int_to_arabic:n { \l_week_number_W_int } }
  }
}

\ExplSyntaxOff


% код для ключа image. ключ имеет 2 аргумента --- ширина и название рисунка
% картинка вставляется в ячейку календаря

\pgfkeys{/tikz/image/.code 2 args=
 {
  \draw (-\cellwidth,0) % задали точку 
  node [above right]    % разместили объект над точкой и правее
   {
   \begin{tikzpicture}
     \includegraphics[width=#1]{#2} 
   \end{tikzpicture}
   };
  }
}


\pgfkeys{/tikz/event/.code =
 {
    \addtext{#1}
  }
}


\pgfkeys{/tikz/anniversary/.code 2 args=\addtextyear{#1}{#2}}

% описываем вид ячейки для даты (прямоугольник со скруглёнными краями)
\pgfkeys{/tikz/day code =
  {
    \node (lower right) at (0,0) [above left,font=\Huge] {\tikzdaytext};
%    \node (upper right) at (-\cellwidth,\cellheight)
%    [below right,align=left,text width=\cellwidth-\pgflinewidth,
%    font=\tiny,black] {\showtext};
        \node (upper left) at (-\cellwidth,\cellheight)
    [below right,align=left,text width=\cellwidth-\pgflinewidth,
    font=\tiny,black] {\showtext};
    \node (lower left) at (-\cellwidth,0) {};
    \node[rounded corners, draw,
          fit=(lower right) (upper left) (lower left),
          inner sep=1mm] {};
  }
}



\pgfkeys{/tikz/inner sep = 0pt}

\pgfkeys{/tikz/day xshift=\cellwidth+3mm+4mm}% xshift --- горизонтальный сдвиг между датами

\pgfkeys{/tikz/day yshift=\cellheight+2mm+2mm}% yshift --- вертикальный сдвиг между датами

\newlength{\cellheight}
\setlength{\cellheight}{18mm}
\newlength{\cellwidth}
\setlength{\cellwidth}{34mm}

% настройка колонтитулов

%\usepackage{fancyhdr}
%\pagestyle{fancy}
\definecolor{navy}{rgb}{0.000000,0.000000,0.501961}


%\renewcommand*{\chapterpagestyle}{empty}

\pagestyle{scrheadings} % активируем команды для задания колонтитулов, описаннные ниже
%\ohead{\pagemark} % номер страницы на чётных страницах слева, на нечётных --- справа

\lefoot{}                      %% Bottom left on even pages
\lofoot{}                      %% Bottom left on odd pages
\refoot{}                      %% Bottom right on even pages
%\rofoot{\color{Black}\sf\small{\copyright{} НТЦ Метротек, 2014}} %% Bottom right on odd pages
\rofoot{\color{Black}\sf\small{}} %% Bottom right on odd pages

\cfoot{\color{navy}\footnotesize\sf{НТЦ Метротек, +7 (812) 340-0118, 
+7 (812) 340-0119 \\ metrotek.spb.ru}}                           %% Bottom center

\lehead{\bfseries\pagemark}    %% Top left on even pages
\lohead{\color{navy}\sf\textbf{\\ \hspace{2mm} \\ \includegraphics[width=45mm]{pics/2014}
}}    %% Top left on odd pages
\rehead{}    %% Top right on even pages
\rohead{\color{navy}\sf\textbf{ }}    %% Top right on odd pages
%\chead{}                           %% Top center

\setkomafont{pagehead}{%
\normalfont\normalcolor\sf\large %шрифт для колонтитулов
}

%%%% высота верхнего колонтитула (header)
\setlength{\headheight}{10mm}

%%%% вертикальный интервал от базовой линии нижней строки текста до
%%%% базовой линии нижнего колонтитула (footer)
\setlength{\footskip}{10mm}

%\deftripstyle{my}{\includegraphics[width=48mm]{pics/2015}}{}{}{}{\color{navy}\footnotesize\sf{НТЦ Метротек, +7 (812) 340-0118, +7 (812) 340-0119 }}{}

\deftripstyle{y2014}{\includegraphics[width=48mm]{pics/2014}}{}{}{}{\raisebox{1ex}{\begin{tabular}[t]{cc} \color{navy}\footnotesize\sf{НТЦ Метротек, metrotek.spb.ru}\end{tabular}}}{}
\deftripstyle{y2015}{\includegraphics[width=48mm]{pics/2015}}{}{}{}{\raisebox{1ex}{\begin{tabular}[t]{cc} \color{navy}\footnotesize\sf{НТЦ Метротек, metrotek.spb.ru}\end{tabular}}}{}


%\setheadtopline{.1pt}[\color{blue}]
%\setheadsepline{}

%\let\endtitlepage\relax




